\section{A law-of-the-iterated-logarithm envelope for the holdover leg}
\label{sec:lil}

The holdover leg raises a foundational question: is there a finite deterministic
number that bounds the coast error for all time? Baweja's impossibility result
\citep{baweja2026tpl} says no, for the single-source case. The boundary-crossing
behaviour of clock-error processes is itself a subject of metrology---the
first-passage / boundary-crossing treatment of Galleani, Tavella, and Peskir
\citep{galleani2016}---which the present LIL application sharpens into an
integrity-risk statement about the holdover envelope's growth order. Here we sharpen
Baweja's result into a two-part statement---an impossibility and a matching
consistency---using the law of the iterated logarithm (LIL). Both parts are \Proven{}, and both are
\emph{asymptotic}: this is stated at the outset because it bounds what the section
may be read to certify.

\subsection{Impossibility of a finite worst-case holdover}
\label{sec:impossible}

The white-FM phase error is, in the limit, a Brownian motion with diffusion $D$.
The Hartman--Wintner law of the iterated logarithm \citep{hartman1941} gives its
almost-sure envelope,
\begin{equation}
  \limsup_{t\to\infty}\frac{|W(t)|}{\sqrt{2Dt\,\ln\ln t}} \;=\; 1
  \quad\text{a.s.},
  \label{eq:hw}
\end{equation}
so $\sqrt{2Dt\,\ln\ln t}\to\infty$: the sample path exceeds any fixed level
infinitely often.

\begin{proposition}[No finite deterministic worst-case holdover bound]
\label{prop:impossible}
For a white-FM (Brownian) holdover leg with $D>0$, there is no finite constant $B$
with $\Prob\big(\sup_{t\ge0}|W(t)|\le B\big)=1$. Equivalently, the almost-sure
envelope \eqref{eq:hw} diverges, so any deterministic ``worst-case holdover''
number is exceeded almost surely.
\end{proposition}

This is why the holdover PL must be \emph{probabilistic}---an integrity-risk-allocated
envelope---and never a worst-case bound. The engine's test
\code{envelope\_diverges\_no\_finite\_worst\_case} checks that the numerical LIL
envelope grows without bound over increasing horizons. The result is
component-specific: Hartman--Wintner \eqref{eq:hw} governs the white-FM/Brownian
leg. The random-walk-FM and drift legs are \emph{integrated} Brownian motions, whose
suprema also diverge almost surely; the relevant statement is the \emph{upper-class}
law-of-the-iterated-logarithm growth (the Hartman--Wintner \citep{hartman1941}
flavour), not the \emph{lower-class} liminf law of Chung \citep{chung1948}---so the
divergence conclusion holds for each leg.

\subsection{Consistency of the integrity-risk-allocated envelope}
\label{sec:consistency}

Impossibility of a worst-case bound does not condemn the probabilistic envelope; it
constrains it. The complementary result is that the IR-allocated $\HPL$ leg grows
\emph{at least as fast} as the a.s.\ envelope once the integrity risk is made
tight.

\begin{proposition}[Consistency as $\IR\to0$]
\label{prop:consistency}
For the white-FM leg, the IR-allocated holdover term
$\Kf(\IR/2)\sqrt{Dt}$ dominates the Hartman--Wintner envelope
$\sqrt{2Dt\ln\ln t}$ precisely when
$\Kf(\IR/2)\ge\sqrt{2\ln\ln t}$. Since $\Kf(\IR/2)=\Phi^{-1}(1-\IR/4)\to\infty$ as
$\IR\to0$, for any fixed horizon $t$ there is an $\IR(t)>0$ below which the
allocated envelope dominates the a.s.\ path envelope; a coupled schedule
$\IR(t)\to0$ keeps the $\HPL$ leg at or above the a.s.\ envelope as $t\to\infty$.
\end{proposition}

The engine's test \code{tightening\_ir\_makes\_hpl\_dominate\_the\_as\_envelope}
exhibits both regimes at a fixed horizon: a loose $\IR=10^{-1}$ leaves the $\HPL$
leg below the a.s.\ envelope, while a tight $\IR=\num{e-9}$ lifts it above.
Proposition~\ref{prop:consistency} is thus the positive counterpart to
Proposition~\ref{prop:impossible}: no finite bound exists, but the allocated
envelope has the \emph{right growth order} to overbound the almost-sure path once
the risk allocation is tight enough.

\subsection{What the LIL certifies, and what it does not}
\label{sec:lilscope}

We state the reach of this section sharply, because the LIL is easy to over-read.
It is an \emph{asymptotic} ($t\to\infty$, $\limsup$) statement about a
\emph{single} sample path. It says nothing about the finite-$t$ distribution and,
in particular, it does \emph{not} certify the finite-$t$ \num{e-7} integrity
quantile: that quantile remains \Modelled{}. The LIL is used here only to
establish two qualitative facts---that no finite worst-case number exists, and
that the allocated envelope's growth order is not sub-envelope---and the
``consistency'' of Proposition~\ref{prop:consistency} is with respect to the
stated coupled $\IR(t)\to0$ schedule, not a claim about any fixed operating point.
Taken together with Section~\ref{sec:impossible}, this turns the known
impossibility into a positive, cited theorem about growth order, while leaving the
finite-horizon coverage question to be answered empirically---which is the subject
of Section~\ref{sec:validation}.
